% Originally: content/week-08.tex, \section{Feynman's trick} (Corsin Week 8)

% Source: Corsin Nick/Class Notes/Week 8.pdf, p. 11
\section{Feynman's trick}
\label{sec:feynmans_trick}

\begin{theorem}[Differentiation under the integral sign]
\label{thm:feynmans_trick}
Let $U \subseteq \mathbb{R}^n\times\mathbb{R}$ be open, $f \in C^1(U)$, and let
$K \subseteq \mathbb{R}^n$ be compact such that $K\times[a,b] \subseteq U$. Then the function
\[(a,b) \to \mathbb{R}, \qquad y \mapsto \int_K f(x,y)\,dx\]
is continuously differentiable, and
\[\frac{d}{dy}\int_K f(x,y)\,dx = \int_K \frac{\partial}{\partial y}f(x,y)\,dx
\qquad \text{for all } y \in (a,b).\]
\end{theorem}

% Source: Corsin Nick/Class Notes/Week 8.pdf, pp. 11--13
\begin{example}[$\int_0^1 \frac{\log(1+x)}{1+x^2}\,dx$ by Feynman's trick]
\label{ex:feynman_log_integral}
We want to compute the integral
\[I(1) = \int_0^1 \frac{\log(1+x)}{1+x^2}\,dx,\]
which has no convenient antiderivative. By introducing a parameter $\alpha \in (-1,1)$:
\[I(\alpha) := \int_0^1 \frac{\log(1+\alpha x)}{1+x^2}\,dx.\]
Note that the integrand is a $C^1$ function on
$U := \big(-\tfrac{1}{1+\varepsilon},\ \tfrac1\varepsilon\big) \times (-\varepsilon,\ 1+\varepsilon)$
in the variables $(x,\alpha)$, with $K := [0,1]$ and $[a,b] := [0,1]$ satisfying
$K\times[a,b] \subseteq U$ for $0 < \varepsilon < 1$ (pay attention to the logarithm, which is
what forces $1+\alpha x > 0$). So, we may differentiate under the integral!

\textbf{We obtain an ordinary differential equation.} By \cref{thm:feynmans_trick},
\[I'(\alpha) = \frac{d}{d\alpha}\int_0^1 \frac{\log(1+\alpha x)}{1+x^2}\,dx
  = \int_0^1 \frac{x}{(1+\alpha x)(1+x^2)}\,dx.\]
Partial fractions: writing
\[\frac{x}{(1+\alpha x)(1+x^2)} = \frac{A}{1+\alpha x} + \frac{Bx+C}{1+x^2}\]
and clearing denominators gives $A + Ax^2 + Bx + B\alpha x^2 + C + C\alpha x = x$, so comparing
coefficients,
\[-A = C, \qquad -\tfrac{A}{\alpha} = B, \qquad B + C\alpha = 1
  \implies -\tfrac{A}{\alpha} - A\alpha = 1,\]
whence
\[A = -\frac{\alpha}{1+\alpha^2}, \qquad B = \frac{1}{1+\alpha^2}, \qquad C = \frac{\alpha}{1+\alpha^2}.\]
Therefore
\[
\begin{aligned}
I'(\alpha) &= \left(\int_0^1 \frac{-\alpha}{1+\alpha x} + \frac{x+\alpha}{1+x^2}\,dx\right)\frac{1}{1+\alpha^2} \\
&= \frac{1}{1+\alpha^2}\Big[-\log(1+\alpha x) + \tfrac12\log(1+x^2) + \alpha\arctan(x)\Big]_0^1 \\
&= \frac{1}{1+\alpha^2}\Big[\log\sqrt2 + \tfrac{\alpha\pi}{4} - \log(1+\alpha)\Big].
\end{aligned}
\]

\textbf{With initial value} $I(0) = \int_0^1 \frac{\log(1)}{1+x^2}\,dx = 0$, we obtain from the
fundamental theorem of calculus
\[
\begin{aligned}
\int_0^1 \frac{\log(1+x)}{1+x^2}\,dx = I(1) &= I(1)-I(0) \\
&= \int_0^1\left(\frac{\log\sqrt2}{1+\alpha^2} + \frac{\pi}{4}\,\frac{\alpha}{1+\alpha^2}
   - \underbrace{\frac{\log(1+\alpha)}{1+\alpha^2}}_{\textstyle \to\, I(1)}\right)d\alpha.
\end{aligned}
\]
The last term is $I(1)$ again, the integral we started from with $\alpha$ merely renamed, so
bringing it to the left-hand side,
\[
\begin{aligned}
2\,I(1) &= \Big[\log\sqrt2\,\arctan(\alpha) + \tfrac{\pi}{8}\log(1+\alpha^2)\Big]_0^1
 = \frac{\pi\log\sqrt2}{4} + \frac{\pi\log\sqrt2}{4}, \\[1ex]
\implies\quad I(1) &= \frac{\pi\log\sqrt2}{4} = \boxed{\ \frac{\pi\log 2}{8}\ }.
\end{aligned}
\]
\end{example}

% Source: Jérôme Paschoud/Notizen/Woche 9.1 Variablenwechsel.pdf, p. 4
% Generator: Gemini 3.1 Pro (High)
\begin{example}[Dirichlet integral via Feynman's trick]
\label{ex:feynman_dirichlet_integral}
Consider the integral $\int_0^\infty \frac{\sin x}{x} \,dx$. We can evaluate it by introducing a damping factor $e^{-tx}$ and defining
\[I(t) := \int_0^\infty e^{-tx} \frac{\sin x}{x} \,dx \qquad \text{for } t > 0.\]
Differentiating under the integral sign with respect to $t$ (which is justified since the integrand and its derivative decay rapidly for $t>0$):
\[I'(t) = \int_0^\infty \frac{\partial}{\partial t} \left( e^{-tx} \frac{\sin x}{x} \right) \,dx = \int_0^\infty -x e^{-tx} \frac{\sin x}{x} \,dx = -\int_0^\infty e^{-tx} \sin x \,dx.\]
This standard integral can be evaluated using integration by parts twice, or by using Euler's formula $e^{ix} = \cos x + i\sin x$ and noting that $e^{-tx}\sin x = \operatorname{Im}(e^{(-t+i)x})$:
\[\int_0^\infty e^{(-t+i)x} \,dx = \left[ \frac{e^{(-t+i)x}}{-t+i} \right]_0^\infty = 0 - \frac{1}{-t+i} = \frac{t+i}{t^2+1}.\]
Taking the imaginary part gives $\int_0^\infty e^{-tx} \sin x \,dx = \frac{1}{t^2+1}$. Thus,
\[I'(t) = -\frac{1}{t^2+1} \implies I(t) = -\arctan(t) + C.\]
Since $\lim_{t \to \infty} I(t) = 0$, we must have $C = \lim_{t \to \infty} \arctan(t) = \frac{\pi}{2}$.
Therefore, $I(t) = \frac{\pi}{2} - \arctan(t)$.
The original integral corresponds to the limit as $t \to 0^+$ (which requires an extra justification step usually provided by Abel's theorem, but we formally take the limit):
\[\int_0^\infty \frac{\sin x}{x} \,dx = \lim_{t \to 0^+} \left( \frac{\pi}{2} - \arctan(t) \right) = \frac{\pi}{2}.\]
\end{example}

\begin{ainote}
Corsin's final integration writes the $\alpha$-term as
$\tfrac{\pi}{2}\cdot\tfrac{\alpha}{1+\alpha^2}$, integrating to $\tfrac{\pi}{4}\log(1+\alpha^2)$.
Both are a factor $2$ away from what $I'(\alpha)$ gives: his own line above has
$\tfrac{\alpha\pi}{4}$, so the term is $\tfrac{\pi}{4}\cdot\tfrac{\alpha}{1+\alpha^2}$ and its
antiderivative is $\tfrac{\pi}{8}\log(1+\alpha^2)$. Corrected above, and the correction is
self-checking, since it is exactly what makes his two bracket terms come out equal, as his own
next line asserts:
$\tfrac{\pi}{8}\log 2 = \tfrac{\pi}{8}\cdot 2\log\sqrt2 = \tfrac{\pi\log\sqrt2}{4}$. His final
answer $\tfrac{\pi\log 2}{8}$ is correct and matches the known value.
\end{ainote}

% Generator: Claude Opus 5 (Medium)
\begin{ainote}[Why this is called a trick, and when to reach for it]
The move that makes this work is not the differentiation. It is noticing that the
\emph{antiderivative} of the resulting integrand is elementary while the original one is not.
The pattern is worth naming:
\begin{enumerate}[label=\textbf{(\alph*)}]
  \item Insert a parameter $\alpha$ so that the target is $I(\alpha_1)$ for some
    $\alpha_1$, and so that $I(\alpha_0)$ is trivial for some other $\alpha_0$. Here
    $I(0)=0$, because $\log 1 = 0$.
  \item Differentiate under the integral. The $\log$ disappears, leaving a rational function
    that partial fractions can handle.
  \item Integrate $I'$ back from $\alpha_0$ to $\alpha_1$.
\end{enumerate}
The self-reference in step 3, where the answer reappears on the right-hand side so that one
solves for $I(1)$ rather than computing it, is specific to this integral and is the reason it
is a classic. Do not expect it in general. Usually step 3 is an honest integration.
\end{ainote}

\begin{aiexercise}[Frullani's integral]
% Generator: Claude Opus 5 (High)
\label{ex:ai_frullani}
Fix $0 < a < b$. The aim is to show that
\[\int_0^\infty \frac{e^{-ax} - e^{-bx}}{x}\,dx = \log\!\left(\frac{b}{a}\right),\]
an answer that depends on $a$ and $b$ only through their ratio.
\begin{enumerate}[label=\textbf{(\alph*)}]
  \item The integrand is undefined at $x = 0$. Compute $\lim_{x\to0^+}$ of it, so that the
    integrand extends continuously to $[0,\infty)$.
  \item Hold $a$ fixed and let the second exponent be the parameter: define
    \[I(t) := \int_0^\infty \frac{e^{-ax} - e^{-tx}}{x}\,dx \qquad \text{for } t > 0.\]
    Differentiate under the integral sign and evaluate the resulting integral, which is
    elementary.
  \item Identify the value of $t$ for which $I(t)$ is trivially zero, and integrate $I'$ from
    there to $t = b$.
\end{enumerate}
\end{aiexercise}
